\def\topic{The free heat equation and uniqueness}
\ifdefined\MyHeaderIncluded
\else
\def\qwertyflag{true}

\ifdefined\MyHeaderIncluded

\else
  \newcommand{\foobar}{FooBar}
  \def\MyHeaderIncluded{true}
% if is till the end of the file

\documentclass[10pt,a4paper]{article}
\usepackage[utf8]{inputenc}
%\usepackage[german]{babel}
%\usepackage[cyrillic]{babel}
\usepackage[T1]{fontenc}
%\usepackage[T2A]{fontenc}
%\usepackage[english, russian]{babel}
\usepackage{amsmath}
\usepackage{amsfonts}
\usepackage{amssymb}
\usepackage{amsthm}
\usepackage{mathtools} %coloneqq

\usepackage{bussproofs} %

%\renewcommand{\thesubsection}{\arabic{subsection}}
%\newtheorem{thm}{Theorem}[section]
%\renewcommand{\thethm}{\arabic{subsection}.\arabic{thm}}

\usepackage{thmtools}
\usepackage{hyperref}
\hypersetup{
    colorlinks=true,
    linkcolor=blue,
    filecolor=magenta,      
    urlcolor=cyan,
    pdftitle={Overleaf Example},
    %pdfpagemode=FullScreen,
    }
%\usepackage{stix-type1}%\turnediota
\theoremstyle{definition}

\newtheorem{constructor}{Constructor}[section]
\newtheorem{iaxiom}{Axiom}
\newtheorem{axiom}{Axiom}
\newtheorem*{axiom*}{Axiom}
\newtheorem{definition}{Definition}[section]
\newtheorem*{definition*}{Definition}
\newtheorem{remark}{Remark}[section]
\newtheorem*{remark*}{Remark}
\newtheorem{recall}{Recall}[section]
\newtheorem*{recall*}{Recall}

\newtheorem{lemma}{Lemma}[section]
\newtheorem{theorem}{Theorem}[section]
\newtheorem*{theorem*}{Theorem}
\newtheorem{note}{Note}[section]
\usepackage{etoolbox}

%\newcommand*{\@Dashv}[2]{%
%  \reflectbox{$\m@th#1#2$}%
%}
\newcommand*{\separation}[3]{% 1)binder, 2)set, 3)condition
  \{#1\in #2: #3\}
}

% https://www.freecodecamp.org/news/simple-regex-tricks-for-beginners-3acb3fa257cb/
%  \\mathcal\{T\} |---> \topology
%  \\mathcal\{([^\}]*)\}   |--->  $1
\newcommand{\AND}{\ \ \&\ \ }
\newcommand{\any}{\mathrm{\ any}}
\newcommand{\borel}{\mathcal{B}}
\newcommand{\can}{\mathrm{can}}
\newcommand{\canonize}{\mathrm{canonize}}
\newcommand{\closure}{\mbox{cl}}
\newcommand{\complex}{\mathbb{C}}
\newcommand{\converges}{\longrightarrow}
\newcommand{\countable}{\mbox{\ countable}}
\newcommand{\ldefined}{\mbox{\ defined}}
\newcommand{\dom}{\mbox{dom}}
\newcommand{\Dom}{\mbox{Dom}}
\newcommand{\Done}{\noindent \textbf{Done.}\\}
\newcommand{\lexists}{\mbox{\ exists}}
\newcommand{\lforall}{\mbox{\ forall}}
\newcommand{\lunique}{\mbox{\ unique}}
\newcommand{\family}{\mathrm{\ family}}
\newcommand{\Fm}{\mathrm{Fm}}
\newcommand{\FuncSymb}{\mathrm{FuncSymb}}
\newcommand{\function}{\mathrm{\ function}}
\newcommand{\functional}{\mbox{\ functional}}
\newcommand{\IFF}{\ \Longleftrightarrow\ }
\renewcommand{\iff}{\ \longleftrightarrow\ }
\newcommand{\IMPLIES}{\ \Longrightarrow\ }
\renewcommand{\implies}{\ \longrightarrow\ }
\newcommand{\Inf}{\mathrm{Inf}}
\newcommand{\injective}{\mbox{\ injective}}
\newcommand{\injection}{\mbox{\ injection}}
\newcommand{\interior}{\mbox{int}}
\newcommand{\intersection}{\bigcap}
\newcommand{\isempty}{\mathrm{\ empty}}
\newcommand{\LB}{\mathrm{LB}}
\newcommand{\nghbset}{\mathcal{N}}
\newcommand{\nonempty}{\mathrm{\ nonempty}}
\newcommand{\Ord}{\mathrm{Ord}}
\newcommand{\power}{\mathcal{P}}
\newcommand{\pr}{\mathrm{pr}}
\newcommand{\fst}{\mathrm{fst}}
\newcommand{\snd}{\mathrm{snd}}
\newcommand{\PredSymb}{\mathrm{PredSymb}}
\newcommand{\Proof}{\noindent \textbf{Proof.}\\}
\newcommand{\ran}{\mathrm{ran}}
\newcommand{\Ran}{\mathrm{Ran}}
\newcommand{\real}{\mathbb{R}}
\newcommand{\relation}{\mathrm{\ relation}}
\newcommand{\set}{\mathrm{\ set}}
\newcommand{\singleton}{\mathrm{\ singleton}}
\newcommand{\semantics}{\mathrm{semantics}}
\newcommand{\semialgebra}{\mathcal{S}}
\newcommand{\Sup}{\mathrm{Sup}}
\newcommand{\surjective}{\mbox{\ surjective\ on\ }}
\newcommand{\Term}{\mathrm{Term}}
\newcommand{\Type}{\mathrm{Type}}
\newcommand{\topology}{\mathcal{T}}
\newcommand{\total}{\mbox{\ total\ on\ }}
\newcommand{\transitive}{\mathrm{\ transitive}}
\newcommand{\union}{\bigcup}
\newcommand{\lTrue}{\mathrm{True}}
\newcommand{\lFalse}{\mathrm{False}}
\newcommand{\lUncertain}{\mathrm{Uncertain}}
\newcommand{\lif}{\mathrm{if}}
\newcommand{\lthen}{\mathrm{then}}
\newcommand{\lelse}{\mathrm{else}}
\newcommand{\lexception}{\mathrm{exception}}
\newcommand{\lundef}{\mathrm{undef}}
\newcommand{\univocal}{\mathrm{\ univocal}}
\newcommand{\univocalat}{\mathrm{\ univocal\ at\ }}
\newcommand{\UB}{\mathrm{UB}}
\newcommand{\Un}{\mathcal{U}}
\newcommand{\Var}{\mathrm{Var}}

\newcommand{\notvarphi}{\phi}
\renewcommand{\phi}{\varphi}

\newcommand{\comment}[1]{\quad\mbox{(#1)}}
% local commands
%\preto\section{\undef\blipo} %etoolbox
\preto\subsection{\undef\blipo} %etoolbox
% https://tex.stackexchange.com/questions/105684/limit-macro-newcommand-scope-to-section
\usepackage{graphicx}
\usepackage{authblk}

\setlength\parindent{0pt} % whole doc of \noindent
\usepackage{soul}
%\usepackage{mathgl}

\fi % guard
\ifdefined\topic
\else
\def\topic{No Titel}
\fi
\title{\topic}
\author[1]{Georgy Dunaev}
\begin{document}
\maketitle
\fi
\section{The free heat equation and uniqueness}
\label{sec:free_heat_equation}
Work on $\mathbb R^d$, $d\geq1$, with Lebesgue measure. In this
chapter write $L=\sum_{j=1}^d\partial_j^2$ for the differential
Laplacian. Thus the nonnegative operator used in the next chapters is
$H_0=-L$, and its heat equation is $u_t=Lu$.
We use elementary multivariable differentiation, the Gaussian integral,
and dominated convergence. The point of the uniqueness argument is its
explicit growth class: a whole-space initial-value problem does not
acquire unrestricted uniqueness from a bounded-domain maximum principle.
Jost's Theorem~5.1.2 is precisely a maximum principle with such a
growth restriction, while his Section~5.2 constructs the fundamental
solution \parencite{Jost2013PDE}.

\subsection{The Gaussian and its convolution law}
For $t>0$ put
\[
p_t(x)=(4\pi t)^{-d/2}\exp\bigl(-|x|^2/(4t)\bigr).
\]
\begin{lemma}[Elementary Gaussian identities]
\label{lem:free_gaussian_convolution}
The function $p_t$ is positive, has integral one, satisfies
$\partial_tp_t=Lp_t$, and obeys
\[
\int_{\mathbb R^d}p_s(x-z)p_t(z-y)\,dz=p_{s+t}(x-y)
\qquad(s,t>0).
\]
\end{lemma}
\Proof
The substitution $x=2\sqrt t\,q$ and the product formula for the
one-dimensional Gaussian integral give $\int p_t=1$. Direct
coordinate differentiation gives
\[
\partial_tp_t(x)=\left(-\frac d{2t}+\frac{|x|^2}{4t^2}\right)p_t(x),
\qquad
\partial_j^2p_t(x)=\left(-\frac1{2t}+\frac{x_j^2}{4t^2}\right)p_t(x),
\]
so summing the second identity proves the differential equation.
For convolution complete the square:
\[
\frac{|x-z|^2}{s}+\frac{|z-y|^2}{t}
=\frac{s+t}{st}\left|z-\frac{tx+sy}{s+t}\right|^2
 +\frac{|x-y|^2}{s+t}.
\]
The remaining Gaussian integral is
$(4\pi st/(s+t))^{d/2}$. Multiplying by the two normalization
factors yields the asserted formula. All these integrals are absolutely
convergent; their positive integrands also permit Tonelli's theorem.
\Done

\begin{lemma}[The Gaussian Fourier multiplier]
\label{lem:gaussian_fourier_multiplier}
For the unitary Fourier normalization
$\widehat f(\xi)=(2\pi)^{-d/2}\int e^{-ix\cdot\xi}f(x)\,dx$,
\[
\widehat{p_t}(\xi)=(2\pi)^{-d/2}e^{-t|\xi|^2}.
\]
Consequently convolution by $p_t$ has Fourier multiplier
$e^{-t|\xi|^2}$.
\end{lemma}
\Proof
In one dimension let
$F(a)=\int_{\mathbb R}e^{-x^2/(4t)}e^{-iax}\,dx$.
The integrable function $|x|e^{-x^2/(4t)}$ permits differentiation
under the integral. Since
$x e^{-x^2/(4t)}=-2t\,\frac{d}{dx}e^{-x^2/(4t)}$,
integration by parts, with vanishing Gaussian boundary terms, gives
\[
F'(a)=2it\int_{\mathbb R}
 \frac{d}{dx}\bigl(e^{-x^2/(4t)}\bigr)e^{-iax}\,dx
 =-2taF(a).
\]
Together with $F(0)=\sqrt{4\pi t}$, this scalar differential
equation gives $F(a)=\sqrt{4\pi t}\,e^{-ta^2}$.
Taking the product in $d$ coordinates proves the displayed transform.
For an integrable function $f$, Fubini applied to the convolution
integral gives
$\widehat{p_t*f}=(2\pi)^{d/2}\widehat{p_t}\,\widehat f$,
which proves the multiplier assertion on integrable functions.
For $f\in L^2$, Jensen's inequality for the probability density
$p_t$ and Tonelli give
\[
\|p_t*f\|_2^2
\leq\int_{\mathbb R^d}\int_{\mathbb R^d}
          p_t(z)|f(x-z)|^2\,dz\,dx=\|f\|_2^2.
\]
Approximation by $L^1\cap L^2$ functions and Fourier--Plancherel
therefore extend the multiplier identity to $L^2$.
\Done

\begin{theorem}[The solution for bounded continuous initial data]
\label{thm:free_heat_bounded_data}
If $f:\mathbb R^d\to\mathbb R$ is bounded and continuous, then
\[
u(x,t)=\int_{\mathbb R^d}p_t(x-y)f(y)\,dy\quad(t>0),
\qquad u(x,0)=f(x)
\]
is continuous on $\mathbb R^d\times[0,\infty)$, smooth for $t>0$,
and solves $u_t=Lu$. Moreover
$|u(x,t)|\leq\|f\|_\infty$ and $u(\cdot,t)\to f$ locally
uniformly as $t\downarrow0$. The operators $P_tf=u(\cdot,t)$ satisfy
$P_sP_t=P_{s+t}$.
\end{theorem}
\Proof
Absolute convergence and the bound follow from $\int p_t=1$.
Every spatial or time derivative of the Gaussian is a polynomial
in $x-y$ times a Gaussian, with coefficients bounded when $t$ stays
in a compact subinterval of $(0,\infty)$. For $x$ in a fixed compact
set these derivatives admit an integrable majorant in $y$ of the form
$C(1+|y|^N)e^{-c|y|^2}$. Differentiation under the integral is therefore
legitimate to every order;
Lemma~\ref{lem:free_gaussian_convolution} gives the equation.

For a compact set $K$ and a number $\varepsilon>0$, continuity of
$f$ on a compact neighbourhood of $K$ supplies $\delta>0$ such that
$|f(x-z)-f(x)|<\varepsilon$ whenever $x\in K$ and $|z|<\delta$.
Consequently
\[
\sup_{x\in K}|P_tf(x)-f(x)|
\leq\varepsilon+2\|f\|_\infty
                 \int_{|z|\geq\delta}p_t(z)\,dz.
\]
After $z=\sqrt t\,q$, the last integral is a Gaussian tail outside
$|q|\geq\delta/\sqrt t$, and tends to zero. This proves the initial
trace and continuity at $t=0$. Fubini, boundedness of $f$, and the
Gaussian convolution identity prove the semigroup law. For complex
initial data apply the construction to real and imaginary parts.
\Done

\subsection{Why the growth condition gives uniqueness}
A classical subsolution on a time slab $[0,T]$ means a real function
continuous up to $t=0$, with continuous second spatial derivatives
and first time derivative for $t>0$, satisfying $u_t-Lu\leq0$.
At the final time derivatives can be understood from the left.

\begin{lemma}[Maximum principle on a bounded cylinder]
\label{lem:heat_cylinder_maximum}
Let $B_R=\{x:|x|<R\}$ and let $u$ be a classical subsolution on
$\overline B_R\times[0,T]$. Its supremum is at most its supremum
on the parabolic boundary
$(\overline B_R\times\{0\})\cup(\partial B_R\times[0,T])$.
\end{lemma}
\Proof
For $\eta>0$ consider $w=u-\eta t$, which satisfies
$w_t-Lw\leq-\eta$. If $w$ exceeded the parabolic-boundary supremum
of $u$, compactness would give a maximum at $(x_0,t_0)$ with
$x_0\in B_R$ and $t_0>0$. At this point $Lw\leq0$ and
$w_t\geq0$: at $t_0=T$ the latter follows from the left derivative,
and at an interior time it is zero. This contradicts the strict
differential inequality. Thus $u\leq\sup_{\partial_p}u+\eta T$;
letting $\eta\downarrow0$ proves the claim.
\Done

\begin{theorem}[Whole-space maximum principle with Gaussian growth]
\label{thm:heat_gaussian_growth_maximum}
Let $T>0$ be finite and let $u$ be a classical subsolution on
$\mathbb R^d\times[0,T]$. Suppose there are $M,\lambda>0$ with
\[
u(x,t)\leq M e^{\lambda|x|^2}\quad(0\leq t\leq T),
\qquad m=\sup_{x\in\mathbb R^d}u(x,0)<\infty.
\]
Then $u(x,t)\leq m$ throughout the slab.
\end{theorem}
\Proof
First suppose the slab length $S$ satisfies $S<1/(4\lambda)$.
Choose $\tau$ with $S<\tau<1/(4\lambda)$ and define
\[
b(x,t)=(\tau-t)^{-d/2}
       \exp\left(\frac{|x|^2}{4(\tau-t)}\right),\qquad0\leq t\leq S.
\]
Direct differentiation gives $b_t=Lb$. For a fixed $\varepsilon>0$,
$w=u-m-\varepsilon b$ is a subsolution and $w(x,0)\leq0$.
On $|x|=R$ the uniform bound
\[
w(x,t)\leq M e^{\lambda R^2}+|m|
 -\varepsilon\tau^{-d/2}e^{R^2/(4\tau)}
\]
is negative for all sufficiently large $R$, since
$1/(4\tau)>\lambda$. The bounded-cylinder lemma implies $w\leq0$
inside every such cylinder. For each fixed $(x,t)$ choose a cylinder
containing it and then let $\varepsilon\downarrow0$. Hence $u\leq m$
on this short slab.

For an arbitrary finite $T$, partition $[0,T]$ into finitely many
slabs of length less than $1/(4\lambda)$. The same growth constants
work on each translated slab. The first step bounds the initial
values of the next slab by $m$, so induction gives the conclusion
on all of them. No limit as a spatial boundary tends to infinity
has been taken without the barrier estimate.
\Done

\begin{corollary}[Uniqueness in the stated growth class]
\label{cor:heat_gaussian_growth_uniqueness}
Two real classical solutions with the same continuous initial data
coincide if, on every finite time slab, each satisfies an absolute
bound $|u(x,t)|\leq M e^{\lambda|x|^2}$ for some slab-dependent
constants. In particular the bounded solution in
Theorem~\ref{thm:free_heat_bounded_data} is unique among bounded
classical solutions. The same conclusion holds for complex solutions.
\end{corollary}
\Proof
On a fixed slab the difference $w$ has zero initial data and an
absolute Gaussian-growth bound, after increasing the two constants
if necessary. Apply the maximum principle to $w$ and to $-w$.
Both are at most zero, hence $w=0$. For complex solutions apply this
to the real and imaginary parts of the difference.
\Done

The Gaussian here is a kernel of the heat semigroup, not a statement
that the differential operator $L$ has only one distributional kernel
or that all unrestricted classical Cauchy problems are unique.
The $L^2$ interpretation uses Fourier--Plancherel and the spectral
calculus in Section~\ref{sec:bounded_schrodinger}; its initial condition
is an $L^2$ limit, whereas the present theorem assumes continuous
initial data and proves a locally uniform limit.
\ifdefined\qwertyflag
\undef\qwertyflag
\end{document}
\fi
